
\section*{Ventana N74: monodromía de nueve puertas}
Definimos:
\[
g(K)=K\bmod9,\qquad 0\equiv9.
\]
Entonces:
\[
\boxed{g(K+9)=g(K).}
\]
Pero:
\[
\boxed{\mathcal S_{K+9}\neq\mathcal S_K\quad\text{en general}.}
\]
La torre no es periódica: es monodrómica.
\[
\boxed{\text{la raíz vuelve; el borde recuerda.}}
\]
La ambigüedad de frontera es de orientación:
\[
u^\varphi=2(q-a),
\qquad
u^\pi+u^e=q-u^\varphi.
\]
Así:
\[
\boxed{\varphi=\text{eje fijo},\qquad \pi/e=\text{espejo móvil}.}
\]
La supervivencia futura rompe el espejo y selecciona la orientación viva.
La novena puerta \(K=9\) no es stutter: es una puerta de orientación donde la monodromía se hace visible.
